\section{Discussion}\label{sec:discussion}

\subsection{Three counts, and why they belong together}

The two halves of this paper look unrelated. One is about attention kernels, the
other about language models applied to tabular rows. The argument is the same in
both. A substitution justified by an adjective (``sparse'', ``pretrained'') is
replaced by a substitution justified by a count, and the count is checked before
the adjective is trusted. For attention the counts are the global-key budget
$A(n)$ of Definition~\ref{def:budget} and the granularity inflation $B/g$ of
Proposition~\ref{prop:block}; for the language model the count is the number of
labelled rows per decision class. Each is a small integer or a simple ratio,
available before training, and in our experiments each predicts the sign of the
result.

The counts also explain \emph{why} the defaults are attractive, which matters for
anyone trying to change practice. Each default removes something awkward to
obtain, namely a task-specific global-token set, a kernel written by hand for one
pattern, or a labelled dataset, and puts something cheap in its place. Under the
constraint that produced it, the substitution is rational. What the measurements
add is the price list. The awkward quantity is worth something, and the cheap
substitute can be an order of magnitude worse on the device, or in the metric the
deployment actually uses.

\subsection{On reporting refutations}

Three claims of our own were refuted during this work and all three are reported
here: the periodicity of real entity sequences
(Section~\ref{sec:acf}), the block-alignment lemma
(Section~\ref{sec:attn_withdraw}), and the premise that a small locally hosted
model comes within two points of a tabular baseline (Section~\ref{sec:lang}). A
fourth, the claim that a sparse pattern saves device memory and latency, survives
only against a baseline that no real system would use, and we say so instead of
quoting the favourable comparison (Section~\ref{sec:attn_disclose}).

We report them because the alternative is worse for the field. Anecdotally, the
failure mode of applied efficiency work is not fabrication but attrition: the
comparison against the naive baseline is kept, the comparison against the fused one
is dropped, the run that produced a null result is not reported, and the paper that
emerges is internally consistent and externally useless. The protocol of
Section~\ref{sec:audit} is a defence against attrition that does not depend on the
author's discipline at the moment of writing: a hypothesis that came out false is a
first-class object in the result file, and the paper has a place to put it.

\subsection{Practical guidance}\label{sec:guidance}

The following procedure is what we would recommend to a team about to adopt either
default, and it is short enough to be used before a project starts.

\begin{enumerate}\itemsep2pt
\item \textbf{Write the mask's entry count, not its complexity class.} Compute
$\lvert\Sset\rvert$ as a function of $n$ for the lengths you will actually serve.
If the count contains a term $n^2/g$ with a fixed $g$, the pattern is quadratic and
the adjective ``sparse'' is describing a constant factor.
\item \textbf{Compute the global-key budget.} If the number of keys visible to
every query grows with $n$, the pattern cannot be linear. Cap the budget first, and
choose the pattern afterwards.
\item \textbf{Divide by the block size.} The executed density is
$\min(1,B/g)$ plus the window term; multiply by $B/B$ to obtain what the kernel
will actually touch. A pattern that is cheap at element granularity can be
expensive at tile granularity.
\item \textbf{Compare against the fused baseline, and against the no-anchor
variant.} The first tells you whether the pattern beats what a real system runs;
the second tells you whether the anchors you added are worth their cost. In our
measurement the second comparison is the one that decides the question.
\item \textbf{Separate interface from model.} When a language model is used as a
decision maker, measure the parse-free interface and the generative one on the same
rows; the difference is often larger than the difference between models, and only
one of the two has a failure mode.
\item \textbf{Count the rows.} Before adapting, state how many labelled examples
per class are available. If the answer is in the hundreds, expect the adaptation to
change the format and not the boundary, and measure both.
\end{enumerate}

\subsection{Threats to validity}

\paragraph{Internal} Splits are strictly temporal and never shuffled; all
fit-dependent statistics are fitted on the training partition only; the operating
point is selected on validation and applied to test once; every comparison between
systems uses identical rows. The two places where internal validity is weakest are
the field ladder, whose smallest sizes require a minority floor, and the
two-point scale fit, which we decline to extrapolate.

\paragraph{Construct} AUC-PR is the primary ranking metric because the positive
class is rare, and it is prevalence-sensitive, which is why we refuse to read the
\DatasetBaf{} drift as robustness. Verifiability is operationalised as an escape
probability under a random challenge, not as a hash-match rate; privacy is not
measured at all beyond the entropy of the salt. Cost per transaction for the second
stage is a proxy constant, not a measured latency, and is labelled as such wherever
it appears.

\paragraph{External} Two datasets, one of them synthetic, and the real one is a
single competition dataset whose slices are not independent external validation.
The entity key is \texttt{card1} because the mirror's user identifier is entirely
null, so entity grouping is coarser than the source intends. One device, one GPU
architecture, one block size, one window. All measurements are single-run except
where seeds or repetitions are reported, and we report repetitions everywhere the
number is a latency.

\paragraph{Conclusion} Ten independent training seeds, clustered and time-block
bootstraps, Holm correction across the two primary tests, effect sizes and an
a-priori power calculation. Where an interval includes zero we say that the
difference is not established, even when the point estimate is favourable.

\paragraph{Hardware and scope} TEE attestation, proof-carrying inference, a
$7$B-class full-precision baseline and public-network gas are not measured. They
are listed with the command that would unblock them and no estimate is substituted.
The quantitative claims of this paper are therefore scoped to a
\GpuMemGiB\,GiB device and a local EVM. The qualitative claims, the two laws, are
scoped only by their hypotheses, which are stated in the theorems.

\subsection{What this work does not establish}

Nothing here establishes that anchor-based sparsity is useless in general. A
fused implementation that materialises the anchor term once per batch, a hardware
pattern able to address strided key blocks without a gather, or a workload whose
$n$ sits far beyond the crossover could each change the verdict, and
Corollary~\ref{cor:ceiling} says how much room there is: a factor of $g/B$, four
at our settings. Nor is there a scale law for language models. Two sizes are two
points, and the inverse ordering at $600$ rows is as likely to be evaluation noise
as an effect. Finally, the decision systems of Section~\ref{sec:system} are not
offered as close to the best achievable on these datasets. Their feature and row
budgets are deliberately small, and a better tabular model would widen the gap
reported in Section~\ref{sec:lang}.

\subsection{Relevance to the wider AI audience}

We close where Section~\ref{sec:intro} began. The two substitutions studied here
are instances of a general move in applied AI: replace an expensive, structured
component by a cheap, general one, and justify it with a scaling argument. The
move is often right. When it is wrong, it is usually wrong for a reason that can be
counted before the experiment, and the count is the part that does not get
published. Our contribution is to make three of those counts explicit, prove them
where they can be proved, measure them where they cannot, and leave the
contradictions in the record. If the reader takes one thing from the paper, we
would prefer it to be the procedure in Section~\ref{sec:guidance}: write the
count first, and let it choose the architecture.
